% !TEX root = ../main.tex
\section{Introducción}

\subsection{Contexto}
La eficiencia energética residencial es un desafío prioritario en Chile. Según el enunciado del proyecto, m\'as del 60\,\% de las viviendas carece de un aislamiento térmico adecuado, lo que incrementa los costos de calefacción y el impacto ambiental asociado. A esto se suman las nuevas exigencias de la Calificación Energética de Viviendas (CEV), que hacen necesario contar con herramientas que permitan a propietarios, arquitectos y constructoras tomar decisiones informadas sobre las soluciones térmicas de los muros, techos, pisos, ventanas y puertas de una vivienda.\\

La dificultad de estas decisiones radica en que cada elemento constructivo puede
materializarse con múltiples alternativas (combinaciones de material y espesor de aislante) que difieren tanto en su desempeño térmico como en su costo, y que ambos criterios están en tensión: una mejor aislación térmica implica, en general, un mayor costo por metro cuadrado.\\

Este proyecto se inspira en una solución real presentada en la Feria de Software
(FESW) 2025 y desarrollada junto a la oficina de arquitectura Acua Plan (sistema
\textit{Bio Planner}), que permite al usuario definir los materiales de algunos
elementos de su vivienda y delegar al sistema la selección del resto desde un
catálogo técnico, minimizando el costo de inversión.

\subsection{Descripción del Problema}
Una vivienda se compone de un conjunto de elementos constructivos pertenecientes a distintas categorías (muros perimetrales, techumbre, pisos, ventanas y puertas), cada uno con una superficie conocida. Para cada categoría existe un catálogo de soluciones constructivas, y cada alternativa del catálogo se caracteriza por dos atributos:\\
\begin{itemize}
	\item \textbf{transmitancia térmica} $U$, unidad que mide el flujo de calor que atraviesa el elemento por unidad de superficie y de diferencia de temperaturas; un valor menor de $U$ indica un mejor aislamiento y, por lo tanto, menores pérdidas de calor;
	\item \textbf{costo} por metro cuadrado de material utilizado.\\
\end{itemize}

Los elementos de la vivienda se clasifican en dos tipos. Los \textit{elementos
fijos} tienen sus materiales predefinidos por el usuario y el modelo no puede modificarlos. Los \textit{elementos optimizables} (o libres) son aquellos para los cuales el modelo debe seleccionar \textbf{exactamente una} solución constructiva del catálogo.\\

\subsection{Objetivo general}
El problema consiste en determinar qué alternativa asignar a cada elemento libre de modo que se \textbf{minimice el costo total de inversión en materiales}, asegurando que la transmitancia térmica promedio de la vivienda, ponderada por el área de cada elemento, no supere un umbral máximo permitido $U_{\max}$.\\

Nuestro objetivo general es modelar matemáticamente este problema de selección de materiales para eficiencia energética residencial como un problema de optimización, generar instancias aleatorias, reproducibles y técnicamente coherentes que permitan evaluarlo, y resolverlas mediante técnicas completas utilizando el solver MiniZinc, obteniendo la solución óptima de cada instancia y verificando su factibilidad y optimalidad. Asimismo, se analiza cómo los márgenes relativos del máximo permitido afectan la factibilidad, el valor objetivo y el espacio de búsqueda, y se evalúa el desempeño computacional del solver en función de la dimensión de cada instancia (variables binarias, restricciones, tiempo de ejecución y, cuando estén disponibles, nodos explorados y brecha de optimalidad).

\subsection{Supuestos Utilizados}
\begin{enumerate}
  \item La vivienda se representa como un conjunto de elementos constructivos
        (muros perimetrales, techumbre, pisos, ventanas y puertas), cada uno con
        una superficie conocida y constante, y perteneciente a una única categoría.
        
  \item Los elementos fijos no pueden modificarse, pero su transmitancia y su costo se consideran en el balance global; cada elemento libre recibe exactamente una solución constructiva del catálogo de su categoría.
  
  \item El costo considerado es solo el de los materiales por metro cuadrado, expresado en unidades monetarias (u.m.); no se incluyen           instalación, mano de obra, transporte ni disponibilidad de stock.
        
  \item Las dimensiones generales de la vivienda (superficie habitable, altura y
        volumen interior) se reflejan únicamente a través de las áreas de sus
        elementos.
        
  \item Como el enunciado no fija valores numéricos, las áreas, transmitancias,
        costos y $U_{\max}$ son datos sintéticos definidos por el generador, bajo el criterio de que a mejor aislación (menor $U$) corresponde mayor costo.\\
\end{enumerate}

Los demás supuestos se describen en su sección correspondiente: los del modelo
matemático en la Sección II-G y los del generador de instancias en la Sección III.
